\paragraph{A covering condition directly on the prime lists.}
The negative rows can be eliminated from the statement. Take a finite
rooted tree with edges labeled by pairs $(r,A)$, where $r$ is an owner
and $A$ a nonempty set of inputs. At a vertex $x$ write $E_x$ for
its outgoing labels. Siblings have distinct labels, and
$E_y\subseteq E_x$ at every child $y$ of $x$. Let $a_x$ give
the last coalition on the path to $x$ at each owner occurring there.
The available colorings at $x$ are those proper passive colorings
$t$ all of whose nonempty blocks occur as labels $(r,t_r)\in E_x$.
Availability at a head $v$ excludes postings of $v$, as usual.
At a leaf only the empty coloring is available.

\begin{theorem}[A prime guard cover on trees of coalitions]
\label{core:thm:prime-tree-cover}
Fix an observer-agreeing antitone family and a requested positive
$(i,p)$. A core word supplies this mark and respects all prescribed
zeros if and only if every tree as above has a vertex $x$ and an
input $v\in L(a_x)$ for which the prime guards at $v$ cover every
available coloring at $x$. Coverage means that a guard is enabled on
the coloring's observer facet. For a complete family with at least
one positive, testing this condition at its maximal positives
characterizes joint realizability.
\end{theorem}
\begin{proof}
If a tree violates the condition, choose at each vertex and each
eligible input one available coloring missed by its guards. For a
vertex $x$ let $S_x$ be all the resulting negative rows in its
subtree. Every eligible projection has a chosen loss in $S_x$. A shortest
feasible word must be singleton at a row of $S_x$, by the descent in
\cref{core:thm:pointed-guard-cut}. If its owner is $r$ and its block
is $A$, the label $(r,A)$ occurs where that row was chosen and hence
at $x$. The word itself meets the updated register at the
corresponding child, including the required gap when $a_x$ was
empty. It still omits all the child's rows. Induction from the
leaves rules this out.

Conversely start with an inclusion-minimal failing set of minimal
losses for the blank register. At a current minimal failing set
$S$ put all its nonempty blocks in $E_x$. Each eligible input
heads some row of $S$, whose coloring is available and missed.
For an outgoing $(r,A)$ delete the rows with $A\subseteq t_r$.
The old register is feasible on the proper subset left. The updated
register must fail there by the two-child construction. Use a
minimal failing subset for that child and repeat. This terminates
and gives the required tree. In fact a child reached through $(r,A)$
retains no label $(r,B)$ with $A\subseteq B$. Finally take the
union of the separate witnesses for the maximal positives.
\end{proof}

The trees in this condition can be large even for short prime lists.
For a cover crown, the same matching rows recur below many different
sequences of offers.

\begin{proposition}[Size of a tree witnessing noncoverage]
\label{core:prop:prime-tree-size}
For the opposite query of a cover crown of length $m$, every tree
witnessing failure in \cref{core:thm:prime-tree-cover} has at least
$2^m$ vertices, already counting only paths that offer the same owner
throughout. On $n$ inputs at fixed $K$, the worst-case minimum size
of a failing tree has double logarithm $n+o(n)$ to base two. The
lower rate holds for polynomial-size canonical prime lists.
\end{proposition}
\begin{proof}
Use the matching intervals $(L_j,U_j)$, indexed from $1$ to $m$
cyclically, and the opposite query $(L_1,U_m)$. Each matching
interval is a cover, so its only negative coloring has its head at
the sole free input and posts $c$ on $L_j$ and $b$ outside $U_j$.
These are all the negative colorings; more than one observer may
read the same row. The rows in any failing tree's root subtree
must include all $m$: deleting any one gives a patch supplying the
opposite mark (\cref{core:lem:crown-patch}). Its root therefore has
all $m$ outgoing labels $(c,L_j)$.

At a vertex reached using only $c$ offers and ending in $(c,L_j)$
the register asks to win on $L_j$ and lose on $L_1$. For every
$1\le k<j$, its subtree must contain matching row $k$. Otherwise
take the greatest bichain completion of the monotone Boolean function
\[
             f(X)=[\,L_h\subseteq X\text{ for some }k<h\le j\,].
\]
It wins on $L_j$ and loses on $L_1$, since the lower endpoints form
an antichain by \eqref{core:eq:crown-incidence}. It omits the mark at
every matching row except $k$: $f(U_\ell)$ reads membership of
$\ell$ or $\ell+1$ in the cyclic interval $(k,j]$, whereas
$f(L_\ell)$ reads membership of $\ell$ alone. This would realize the
register on all the subtree rows, contradicting the preceding proof.
Nested outgoing labels thus force a child $(c,L_k)$ for every $k<j$.
Every decreasing sequence of indices occurs as a distinct path from
the root, including the empty one. There are $2^m$.

For the upper bound put $d=K(2^n-1)$. In the tree constructed in the
theorem no edge label repeats along a path, so there are at most
$\sum_{h=0}^d d!/(d-h)!$ vertices. Its double logarithm is at most
$n+O_K(\log n)$. \Cref{core:cor:compact-locality-rate} gives the
matching lower rate using polynomial-size prime lists. For all
sufficiently large arities the cover crowns of
\cref{core:thm:interval-mask-locality-rate} give
$m\ge 2^{n-O(\log n)}$ without the prime-size restriction.
\end{proof}

Sharing subtrees changes this bound. Allow a finite rooted acyclic
directed graph, with the register written on each vertex and updated
along each edge. The same local test is exact, by unfolding the graph.
For a crown of length $m$ the least number of vertices in such a
certificate is exactly $(m+1)(m+2)/2$. Indeed, after a first offer $(c,L_j)$,
the preceding proof forces matching row $k$ somewhere below it for
every $k<j$. This forces edges $(b,V\setminus U_k)$ at that vertex
as well. Their children carry $\binom m2$ distinct two-owner
registers, in addition to the $2m$ forced one-owner registers and
the blank root.

For the upper bound take the path from $U_m$ to $L_1$ that reverses
matching rows $m,m-1,\ldots,1$ in order, using the inclusion
$L_j\subseteq U_{j-1}$ between reversals. Keep one vertex for each
contiguous segment of this path's list of reversals, including the
$m+1$ empty segments at its cuts. At a segment use the labels
$(c,L_j)$ and $(b,V\setminus U_j)$ of all its rows. The first edge
keeps the suffix after $j$ and the second keeps the prefix before $j$.
The endpoints $P,N$ of the segment start as $U_m,L_1$; a $c$-edge
replaces $P$ by $L_j$, and a $b$-edge replaces $N$ by $U_j$.
The segment boundaries thus determine both the register and its
eligible set $P\setminus N$, including the sole free head at the
root. Each $v\in P\setminus N$ must leave the running set somewhere
on this path. An inclusion cannot remove it, so it is the free head
of a reversal in the segment. That row supplies a missed coloring
using the vertex's outgoing labels. For an empty segment $P\subseteq N$.
This proves noncoverage everywhere with at most
$(m+1)(m+2)/2$ vertices.

\paragraph{Which lattices of remaining rows occur?}
Keep the rows left uncovered by a sequence of offers, irrespective of
the replies. These subsets of $Z$ form a finite lattice $L$ under
inclusion: concatenating two sequences takes the intersection, and
the empty sequence leaves $Z$. Write $z_r$ for the $r$-coalition of
row $z$. For $T\subseteq Z$ the least reachable set containing $T$ is
\begin{equation}\label{core:eq:offer-closure}
 \overline T=\{z\in Z:\ \forall r\text{ with }z_r\ne\varnothing,
                 \exists y\in T,\ z_r\subseteq y_r\}.
\end{equation}
Indeed, every offer that removes such a $z$ also removes a row of $T$.
For a row failing the test, offer its offending block $z_r$. Doing so
for every failing row leaves exactly $\overline T$. In particular,
allowing unobserved coalitions gives the same lattice. Each membership
$z\in\overline T$ uses at most $K-1$ witnesses in $T$.

For minimal losses this lattice distinguishes exactly the histories
that have different answers to a future head-coverage test. If $z$
survives in one history alone, let $j$ be its head. For every row
$y$ of the other history headed by $j$, some block $y_r$ is not
contained in $z_r$, by minimality. Offering those blocks covers
every remaining loss at $j$ in the other history and leaves $z$
in the first. Histories with identical surviving rows give the same
coverage after every continuation. The accepted coalitions in the
second terminal condition of \cref{core:cor:alternating-transversal}
are recorded separately.

Here is an intrinsic description of these lattices. In a finite
lattice write $J$ for the nonzero join-irreducibles and $M$ for the
nonunit meet-irreducibles. Put $M_j=\{m\in M:j\nleq m\}$.

\begin{theorem}[Owner colorings of the lattice of remaining rows]
\label{core:thm:offer-lattice}
Let $K\ge3$. A finite lattice occurs as the lattice of remaining
minimal losses of a core operation on $K$ owners if and only if there
is a coloring $c:M\to I$ such that, for every $j\in J$,
\[
 c(M_j)\ne I,\qquad
 j\nleq\bigvee\{m\in M_j:c(m)=r\}\quad(r\in I).
\]
Every such lattice occurs with one loss per conditional source and
all its postings on retained inputs.
\end{theorem}
\begin{proof}
In a lattice of remaining rows, each element $x$ is the join of the
closures of its individual rows. Thus for $j\in J$ choose $z$ with
$j=\overline{\{z\}}$: membership of $z$ in any $x$ means $j\le x$.
Every element is an intersection of single-offer survivor sets, so each
$m\in M$ is itself the survivor set of one offer. Color $m$ by the
owner of a chosen such offer. If $m\in M_j$, its offer removes $z$;
hence its color occurs on $z$, which omits an owner. For a nonempty
family of offers at one owner the union of their survivor sets is
itself a survivor set: offer the union of their coalitions. Thus this
union is their lattice join. If every one removes $z$, their join
still omits $z$. The empty join lies strictly below $j$. This proves
both conditions.

Conversely represent $x$ by $J(x)=\{j\in J:j\le x\}$. Make one
head for each $j$ and one retained input $u_m$ for each $m\in M$.
At head $j$ prescribe the single minimal loss posting exactly
$u_m:c(m)$ for $m\in M_j$. This coloring is proper by the first
condition; it is nonempty since the meet-irreducibles separate $j$
from the bottom element. These losses occur together in a core
operation: for each head take its one-source completion on the common
retained inputs (\cref{core:thm:source-completion}) and take their
union. A head absent from a summand has no source there. For the
one-element lattice use a projection with no losses.

An offer can remove a row only if it consists of retained inputs
$u_m$ all posted at its owner, say with nonempty index set $X$. Its
survivors are $\bigcup_{m\in X}J(m)$. This equals $J(\bigvee X)$:
a join-irreducible absent from every $J(m)$ is absent from their join
by the second condition. Ineffective offers leave the top. Every
history therefore leaves a set $J(x)$, and offering $u_m$ alone at
$c(m)$ leaves $J(m)$. Since every $x$ is a meet of meet-irreducibles,
all $J(x)$ occur. The order on them is the original lattice order.
\end{proof}

For example, take a height-two lattice with $d\ge3$ atoms. Two
different atoms join above any third, so the atoms need distinct
colors. Conversely $d$ colors suffice, since $M_j$ omits $j$.
Adjoining a new top above this lattice raises the least number of
owners to $d+1$: the new top is join-irreducible and sees all $d$
colors among the atoms it misses. Giving the old top any one of
those colors and leaving the extra color unused meets both conditions.
Thus even among acyclic realizable families the number of owners
gives a strict hierarchy of lattices. A lattice type by itself cannot
exclude a candidate family: every type arising from proper rows
already occurs among realizable families, by the same construction.

For three and four owners the coloring test is ordinary graph
colorability. Put vertices $J\sqcup M$, join $j$ to each member of
$M_j$, and join $m,m'$ if their join lies above a $j$ missing from
both. A lattice coloring extends to this graph by assigning each
$j$ a missing color. Conversely suppose the graph is properly colored.
The maximal members of $M_j$ have distinct colors: the join of two of
them lies above $j$ by maximality. Consider a color class $X$ in
$M_j$. If its color occurs at a maximal member $m$, every member of
$X$ lies below $m$, by the same edge test. Otherwise at most $K-2$
maximal members occur: both the color of $j$ and the color of $X$
are missing among them. If no maximal member bounds $X$, one can
choose a member of $X$ outside each such bound. Their join lies above
$j$, since any element missing $j$ lies below a maximal member of
$M_j$. For $K\le4$ this contradicts the pair test. The same argument
for general $K$ shows that monochromatic subsets of size at most
$\max(2,K-2)$ suffice in \cref{core:thm:offer-lattice}.

\paragraph{An implicational test and the size of its countermodels.}
There is also a test in ordinary intuitionistic implication. It uses
one atom $d_z$ per listed negative and one atom $w_a$ per register
assignment, including the blank assignment $\varnothing$. For a list
$B$ of rows let $[B]w$ abbreviate the nested implications with
premises $d_z$ for $z\in B$ and conclusion $w$; an empty list gives
$w$. Write $Z_v$ for the rows with head $v$ and $C_e$ for those
covered by a nonblank observed offer $e$. Use the hypotheses
\[
 [Z_v]w_a \quad(v\in L(a)),\qquad
 ([C_e]w_a)\ \to\ ([C_e]w_{a[e]})\ \to\ w_a
       \quad\text{(every $a,e$)},
\]
where $a[e]$ replaces the entry at the offered owner.
These formulas have polynomial total length on an explicit sample at
fixed $K$: there are at most $(1+K|Z|)^K$ assignments.

\begin{proposition}[An implicational criterion for guard realization]
\label{core:prop:guard-sequent}
The hypotheses entail $w_\varnothing$ in intuitionistic logic exactly
when the marked query is feasible. On failure the lattice of remaining
rows, ordered by reverse inclusion, supplies a finite Kripke
countermodel. If every failing subfamily of the negatives has at
least $b$ rows, then every finite Kripke countermodel has width at
least $\lceil b/|V|\rceil$.
\end{proposition}
\begin{proof}
Read a winning tree from its leaves upward. At a leaf use the first
hypothesis for its eligible input; each required $d_z$ has already
been supplied on the path. At an offer abstract the assumptions for
$C_e$ in both continuation proofs and apply the second hypothesis.
This gives a derivation of the root atom with no row assumptions.

On the reverse lattice force $d_z$ at $S$ when $z\notin S$, and
$w_a$ when $a$ is feasible on $S$. Both predicates persist upon
deletion. The first world above $S$ forcing every $d_z$ for $z\in C_e$
is $S\setminus C_e$. Thus $[C_e]w_a$ holds at $S$ exactly when
$a$ is feasible on this smaller set. The two-child construction of
\cref{core:thm:pointed-guard-cut} validates the offer hypothesis;
a fitting projection validates the leaf hypothesis. Failure makes
$w_\varnothing$ false at the initial world $Z$.

In any finite countermodel of width $q$ consider only worlds above
and including the refuting root. For each head $v$ the worlds where
at least one $d_z$ for $z\in Z_v$ is false form a downset.
At each of its maximal worlds retain one such row. At most $q|V|$
rows are retained. Persistence means that a false premise still
exists everywhere in that downset. The leaf hypotheses therefore
still hold: whenever $w_a$ is false, every eligible head has a false
premise at that world. Removing premises from $[C_e]w$ makes this
antecedent of the offer rule harder to satisfy, so all remaining offer
hypotheses still hold. Discard the offers and register assignments
absent from the restricted sample. The same root still refutes the
conclusion, so the retained sample fails, proving the bound.
\end{proof}

The implication test has a succinct form on the canonical primes
alone. Use finite, constant index domains for the inputs, owners,
proper passive colorings, nonblank offers and register assignments;
their indices have polynomial bit length. Let $X_r(A)$ be an atomic
predicate on nonempty coalitions, subject to
$X_r(A)\to X_r(B)$ for $A\subseteq B$.
For a coloring $t$ at $v$ write $P_v(t)$ for the rigid input test
that some listed prime guard is enabled on its observer facet.
Agreement makes the choice of a missing observer immaterial. Put
\[
 Q(v)=\forall t\text{ at }v\quad
 \left(P_v(t)\ \lor\!
       \bigvee_{r:\,t_r\ne\varnothing} X_r(t_r)\right).
\]
Use just the following two universally quantified hypotheses, in
addition to monotonicity:
\[
  Q(v)\to W(a)\quad(v\in L(a)),\qquad
  (X_r(A)\to W(a))\to(X_r(A)\to W(a[r:=A]))\to W(a).
\]
The query succeeds exactly when these hypotheses entail
$W(\varnothing)$. Quantifiers here abbreviate finite conjunctions or
finite case distinctions over the stated domains. Keeping the indices
bound gives polynomial length in the prime presentation; expanding
them can take exponential space.

Indeed, along a successful branch the assumed offers cover every
losing coloring of the terminal input. For each such coloring,
monotonicity gives $X_r(t_r)$ from a covering offer, proving $Q(v)$
by cases. Conversely use the same countermodel with $X_r(A)$ true
at $S$ when no row left in $S$ is covered by $(r,A)$. This predicate
is monotone in $A$ and persistent under deletion. The first extension
making it true deletes exactly those rows. If a loss at $v$ survives,
its own coloring falsifies $Q(v)$, since every nonempty block covers
that very row. If none survives, every losing coloring extends a
minimal loss at $v$ already deleted by some offer. That offer and
all enlargements of its coalition are true at $S$. Thus $Q(v)$
holds exactly when the leaf $v$ is available. The two hypotheses
follow as in the explicit sequent. Every countermodel of the indexed
form also gives one for the explicit sequent on the same frame: set
$d_z=\bigvee_{r:z_r\ne\varnothing}X_r(z_r)$, with the empty
disjunction false. Truth of all minimal-loss premises for $v$ implies
$Q(v)$ by monotonicity. Forcing an offer forces $d_z$ for every row
it covers, giving the explicit offer hypothesis.

\begin{proposition}[Shared certificates and finite countermodels]
\label{core:prop:guard-dag-countermodel}
For a failing indexed query, let $D$ be the least number
of vertices in a graph certifying noncoverage and $F$ the least number
of worlds in a countermodel. For $n=|V|$ one has
\[
                         F\le D\le F(1+nF)^K .
\]
\end{proposition}
\begin{proof}
For the first inequality order the graph by descent from its root.
Choose one missed row for each eligible input at each vertex.
Force $X_r(A)$ at $x$ when $(r,A)$ covers no chosen row at or below
$x$. Make $W(a)$ false at $x$ when some descendant $y$, possibly
$x$, has a register on the same owners containing $a$ coordinatewise.
The eligible inputs of $a$ are then eligible at $y$, whose chosen
rows refute their $Q(v)$ at $x$. To check the offer rule for
$e=(r,A)$, start at this $y$. If $e$ covers no chosen row below it,
the old register is still false there and $X_r(A)$ holds.
Otherwise some such row has a block $B\supseteq A$ at $r$.
Heredity gives an edge $(r,B)$ here. Take this edge and repeat until
$e$ covers no remaining row. The graph is finite and acyclic, so the
process terminates. The final register contains $a[r:=A]$, since
only $r$ was updated. Thus at least one of the two antecedents in
the offer rule fails at $x$. Both predicates persist under descent,
and the blank register fails at the root.

Conversely, use pairs $(x,a)$ where $x$ is a maximal world at
which $W(a)$ is false. Start at such a world for the blank register
above the refuting root. For each input $v$
choose a missed coloring at each maximal world failing $Q(v)$.
Its nonempty blocks have false atoms there, since a coloring
witnessing failure at an extension also witnesses it at the original
world. At $(x,a)$ use as labels the blocks $(r,A)$ of the selected colorings
for which $X_r(A)$ is false at $x$.
Every eligible $v$ at $x$ still has a selected coloring all of whose
blocks are false at $x$, by persistence. For each outgoing label $e$, failure of an
offer antecedent
gives a world above $x$ forcing its atom where the old or updated
register fails. This world is strictly above $x$, so maximality
excludes the old register. Extend to a maximal world failing the
updated register and use that pair as the child. Persistence gives
nested labels; monotonicity removes the whole upset of the accepted
coalition. At frame width $u$ at most $nu$ colorings were selected;
there are at most $1+nu$ register entries per owner, including absence.
This gives at most $F(1+nu)^K$ pairs and hence the bound.
\end{proof}

For the opposite query of a cover crown of length $m$ from
\cref{core:thm:exponential-locality}, the least possible countermodel
height is exactly $m+1$, counting worlds. As in \cref{core:prop:prime-tree-size}, write $z_j$ for
matching row $j$, use $(L_1,U_m)$ as the opposite query, and offer
$(c,L_j)$ in order $m,m-1,\ldots,1$. Each offer covers just $z_j$,
by \eqref{core:eq:crown-incidence}. At the root $d_{z_m}$ would
force $w_\varnothing$: deleting $z_m$ gives a patch supplying the
mark. The offer hypothesis for the first step therefore forces a strict successor
where $d_{z_m}$ holds and the register $c:L_m$ fails.

Consider the subsequent step from $c:L_j$ to $c:L_{j-1}$ for $j>1$.
Once $d_{z_{j-1}}$ also holds, the old register must succeed.
Indeed its task is to win on $L_j$ and lose on $L_1$ at $c$.
The Boolean function $X\mapsto[L_j\subseteq X]$ does this, and its
greatest bichain completion omits every row except possibly $z_j$
and $z_{j-1}$: these are the only upper endpoints containing $L_j$.
The derivation above, starting at this register, thus forces its atom
under the two row assumptions. Consequently $d_{z_{j-1}}$ is still
false at the present world. By the offer hypothesis it becomes true
at a strict successor where the new register fails. This gives $m$
strict steps. The reverse lattice on $m$ rows has height at most
$m+1$, proving equality. Every deletion patch supplies the opposite
query, so the width bound is at least $\lceil m/|V|\rceil$ as well.

The same crowns require $\Theta(m^2)$ worlds in total. Offer
$(c,L_j)$ and then $(b,V\setminus U_k)$ for $1\le k<j\le m$. Each forces a successor as above:
the deletion patch at $j$ forces the blank register under $d_{z_j}$,
and the separator for $(k,j]$ in \cref{core:prop:prime-tree-size}
forces the $c:L_j$ register under $d_{z_k}$. At the resulting world
$d_{z_k},d_{z_j}$ are true and the two-owner register fails.
Every $d_{z_h}$ with $k<h<j$ must be false there: the separator
for $(h,j]$ wins on $L_j$, is zero on $U_k$, and omits every mark
except $h$. Keeping $k\le\lfloor m/2\rfloor<j$, the nearest true
indices on either side of this cut therefore recover $k,j$.
This gives at least $\lfloor m^2/4\rfloor$ distinct worlds.
Pairs with both endpoints strictly increasing give
$\lfloor m/2\rfloor$ incomparable worlds, since each forces a row
false at the other. The interval graph gives a countermodel of
quadratic size by the preceding construction and has width at most
$m+1$: segments with a common left boundary are linearly ordered by
reachability (a prefix edge cuts at any row).
Thus minimum width has order $m$, and minimum height is $m+1$.
The lower bounds hold for the explicit as well as the indexed
sequent, by the same row implications. For the polynomial prime
lists of \cref{core:cor:compact-locality-rate} both dimensions can
be $2^{|V|-o(|V|)}$. At width four of the primes both already grow
exponentially (\cref{core:thm:compact-prime-locality}).
